\section{Evaluation}
Im Rahmen dieses Projekts wurden viele funktionale und technische Anforderungen an das zu entwickelnde Visualisierungs-Framework definiert.
Bei genauer Betrachtung lässt sich sagen, dass die meisten dieser Ziele erfolgreich erreicht wurden.

Ein besonders wichtiges Ziel war, das Framework so zu entwickeln,
dass es direkt in JupyterLab genutzt werden kann. Das hat gut funktioniert:
Die Anwendung lässt sich problemlos in Jupyter Notebooks einbinden, was die Nutzung für Studierende unkompliziert und plattformunabhängig macht.
Besonders praktisch ist dabei, dass die 3D-Visualisierung direkt im Notebook erscheint. Das macht die Bedienung und das Ausprobieren deutlich einfacher und anschaulicher.

Ein weiteres wichtiges Kriterium war die interaktive Bedienbarkeit.
Durch die Verwendung von ipywidgets konnten benutzerfreundliche Bedienelemente wie Schieberegler und Dropdown-Menüs implementiert werden.
Diese ermöglichen eine intuitive Steuerung von Parametern, wie zur Anpassung von Rotationen und erfüllen somit die
Anforderung nach einer interaktiven Auseinandersetzung mit kinematischen Konzepten.

Die Kompatibilität mit SymPy, die symbolische Berechnungen und deren Integration in die Visualisierungen erlaubt,
wurde ebenfalls sichergestellt. Die Funktionen in Util.py können problemlos Matrixobjekte der Sympy Bibliothek entgegennehmen.

Im inhaltlichen Fokus stand die Darstellung und Transformation von Rotationen im dreidimensionalen Raum.
Die Umsetzung verschiedener Darstellungsformen wie Eulerwinkel, Roll-Pitch-Yaw-Winkel, Rotationsmatrizen und Quaternionen wurde erfolgreich realisiert.
Die entsprechenden Umrechnungen zwischen diesen Darstellungen sowie deren visuelle Präsentation im Raum stellen eine große didaktische Stärke des Frameworks dar und decken die Anforderungen vollständig ab.

Auch die Simulation von radgetriebenen Robotern wurde erfolgreich integriert. Mit der Funktion move(robot, vel, steps)
lassen sich einfache Bewegungsabläufe in der Ebene darstellen. Das erweitert das Framework sinnvoll über reine Rotationen hinaus.

Die Verwendung von Poetry zur Verwaltung von Abhängigkeiten sorgt darüber hinaus für eine stabile und konsistente Entwicklungsumgebung,
was insbesondere für die langfristige Nutzung und Weiterentwicklung des Frameworks relevant ist.

Einziger Aspekt, der im aktuellen Projektstatus noch nicht vollständig umgesetzt wurde,
ist die Modellierung stationärer Roboter (z.B. Knickarmroboter) samt Vorwärts- und Rückwärtskinematik.
Allerdings wurde dies von Beginn an als zukünftige Erweiterung vorgesehen und fällt nicht in den unmittelbaren Umfang dieser Projektarbeit.

Es lässt sich feststellen, dass das Framework die definierten Anforderungen in nahezu allen Punkten erfüllt.
Es bietet eine funktional umfangreiche, technisch stabile und didaktisch wertvolle Plattform zur Unterstützung der Lehre im Bereich der robotischen Kinematik.
Die Grundlage für eine spätere Erweiterung um komplexere Robotermodelle wurde ebenfalls gelegt, womit auch die langfristigen Projektziele realistisch erscheinen.